\section{LC-ORF Framework}
\label{sec:lcorf}

This paper proposes LC-ORF, a leakage-controlled operational robustness framework for fraud-detection evaluation. LC-ORF is not a new classifier. It is an audit protocol that measures how fraud-model conclusions change when key evaluation assumptions are perturbed. The framework evaluates six main operational axes: leakage, temporal protocol, prevalence shift, calibration, explanation protocol, and alert-budget behavior. A final explanation-reliability-gap component is then derived from the predictive and explanation-stability artifacts.

\subsection{Audit Units and Scenario Pairs}

Let $\mathcal{D}$ denote the evaluated datasets, $\mathcal{M}$ the model families, $\mathcal{S}$ the random seeds, $\mathcal{A}$ the audit axes, and $\mathcal{K}$ the metrics. An LC-ORF audit unit is indexed by $(d,m,s,a,c_0,c_1,k)$, where $d \in \mathcal{D}$, $m \in \mathcal{M}$, $s \in \mathcal{S}$, $a \in \mathcal{A}$, $k \in \mathcal{K}$, $c_0$ is a reference scenario, and $c_1$ is a perturbed scenario. In this study, $\mathcal{D}$ contains D1\_creditcard, D2\_fraudTest, and D3\_PaySim; $\mathcal{M}$ contains LightGBM and XGBoost; and $\mathcal{S}=\{42,101,202,303,404\}$.

For each metric $k$, LC-ORF defines the scenario delta as
\begin{equation}
\Delta_{a,d,m,s,k}
= Q_k(d,m,s,c_1) - Q_k(d,m,s,c_0),
\label{eq:lcorf-delta}
\end{equation}
where $Q_k$ is the metric functional computed from the true labels, model scores, and the decision threshold. This sign convention is perturbed minus reference. Thus, a positive leakage delta indicates that an unsafe protocol increased the measured metric relative to the strict protocol, while a negative temporal delta indicates that chronological evaluation reduced the metric relative to random stratified evaluation. For lower-is-better metrics such as Brier score, expected calibration error, and cost per 10,000 transactions, the sign is interpreted according to the metric direction.

\subsection{Uncertainty Estimation}

For leakage, temporal, calibration, and explanation-protocol comparisons, LC-ORF estimates uncertainty with $B=1000$ bootstrap replicates. If reference and perturbed predictions have row-level alignment, a paired bootstrap is used. Otherwise, LC-ORF uses stratified unpaired resampling to preserve the fraud and non-fraud label composition in each scenario. For bootstrap replicate $b$, the method computes
\begin{equation}
\Delta^{(b)}_{a,d,m,s,k}
= Q^{(b)}_k(d,m,s,c_1) - Q^{(b)}_k(d,m,s,c_0).
\end{equation}
The interval estimate is the percentile interval
\begin{equation}
CI_{a,d,m,s,k}
= \left[
P_{2.5}\left(\Delta^{(1:B)}_{a,d,m,s,k}\right),
P_{97.5}\left(\Delta^{(1:B)}_{a,d,m,s,k}\right)
\right].
\end{equation}
The final audit profile aggregates across seeds by reporting the mean delta, the standard deviation across seeds, the median lower interval bound, the median upper interval bound, the minimum bootstrap count, and the bootstrap mode.

\subsection{Materiality Labels}

LC-ORF assigns robust, fragile, or inconclusive labels using metric-specific materiality thresholds $T_k$. In this study, $T_{AP}=0.020$, $T_{MCC}=0.050$, $T_{precision}=0.050$, $T_{recall}=0.050$, $T_{Brier}=0.005$, and $T_{ECE}=0.020$. For cost per 10,000 transactions, $T_{cost}=\max(100,0.10|Q_{cost}(c_0)|)$.

A comparison is labeled robust when both interval endpoints lie inside the materiality band:
\begin{equation}
|CI_{\mathrm{low}}| < T_k
\quad \text{and} \quad
|CI_{\mathrm{high}}| < T_k.
\end{equation}
It is labeled fragile when the interval is entirely outside the materiality band in either direction:
\begin{equation}
(CI_{\mathrm{low}} > T_k \land CI_{\mathrm{high}} > T_k)
\quad \text{or} \quad
(CI_{\mathrm{low}} < -T_k \land CI_{\mathrm{high}} < -T_k).
\end{equation}
All remaining cases are labeled inconclusive.

\subsection{Prevalence and Alert-Budget Components}

The prevalence component evaluates fixed models and thresholds under altered fraud-rate composition in the evaluation stream. For each target prevalence, LC-ORF records the difference between the target-prevalence metric and the natural test-stream metric. The final prevalence profile reports the mean delta across target prevalence levels, the worst target prevalence, and the maximum absolute delta.

The alert-budget component ranks transactions by model score and evaluates fixed investigation budgets. The reported quantities are alert precision, alert recall, false alerts, and estimated cost. This component connects score behavior with the operational workload faced by fraud analysts.

\subsection{Explanation Reliability Gap}

The explanation-reliability gap is an artifact-level proxy. Let $\overline{J}^{(top-k)}_{d,m}$ be the mean top-$k$ Jaccard agreement between explanation feature rankings under the compared protocols. LC-ORF defines
\begin{equation}
EII_{d,m}=1-\overline{J}^{(top-k)}_{d,m}.
\end{equation}
Using $k=5$ and $T_{AP}=0.020$, the AP fragility unit is
\begin{equation}
F^{AP}_{d,m}
= \frac{|\Delta^{AP}_{explanation,d,m}|}{T_{AP}}.
\end{equation}
The explanation-reliability gap is
\begin{equation}
ERG_{d,m}=F^{AP}_{d,m}-EII_{d,m}.
\label{eq:erg}
\end{equation}
An $ERG$ above 1 indicates that predictive degradation exceeds explanation-rank instability by at least one AP materiality unit. This quantity should not be interpreted as causal explanation validation. It is a conservative diagnostic that reports disagreement between predictive behavior and explanation-rank movement in the available artifacts.

\subsection{Final Audit Profile}

The final LC-ORF profile joins the bootstrap-backed deltas, prevalence profile, explanation-reliability gap, and alert-budget summaries into an auditable result set. The purpose of this profile is not to declare a single winning model, but to show where fraud-detection conclusions are robust, fragile, or inconclusive under changed operational assumptions.
